\section*{HISTORICAL NOTE}
\phantomsection
\addcontentsline{toc}{section}{Historical Note}

(Numbers in brackets refer to the bibliography at the end of this note.)

We shall not repeat here the complete history of the development of the theories of complex numbers and quaternions, since these theories belong essentially to algebra; but we shall say something about the geometrical representation of complex numbers, which in many respects was a decisive step forward in the history of mathematics.

Without any doubt the first to have a clear conception of the one-to-one correspondence between complex numbers and points in the plane was C. F. Gauss (*), who moreover applied this idea to the theory of complex numbers and foresaw the uses to which it would be put by the analysts of the 19th century. During the 17th and 18th centuries, mathematicians had gradually come to the conviction that imaginary numbers, which allowed them to solve all quadratic equations, could also be used to solve algebraic equations of any degree. Many attempts to prove this theorem were published during the course of the 18th century; but, not mentioning those which were based only on a vicious circle, there was not one which was not open to serious objection. Gauss, after a detailed examination of these attempts and a closely reasoned criticism of their deficiencies, proposed in his inaugural dissertation (written in 1797, published in 1799) to give at last a rigorous proof. Taking up an idea thrown off in passing by d'Alembert {[}in his proof published in 1746 (**){]},

Gauss remarked that the points \((a, b)\) of the plane which are such that \(a + ib\) is a root of the polynomial \(\mathrm{P}(x + iy) = \mathrm{X}(x, y) + i\mathrm{Y}(x, y)\) are the intersections of the curves X = o and Y = o; by means of a qualitative study of these curves, he then shows that a continuous arc of one of them joins points of two distinct regions bounded by the other, and concludes from this that the curves intersect ({[}1{]}, vol.~3, p.~3; see also {[}1 a{]}). In its clarity and originality this proof was a considerable advance on previous attempts, and was certainly one of the first examples of purely topological reasoning applied to a problem of algebra (*).

In his dissertation, Gauss does not explicitly define the correspondence between points of the plane and complex numbers; as to the latter, and the questions of ``existence'' to which they had given rise for two hundred years, he reserves his position and deliberately presents all his arguments in a form which involves only real quantities. But the march of ideas in his proof would be wholly unintelligible if it did not presuppose a fully conscious identification of points of the plane with complex numbers; and his research at the same period in number theory and elliptic functions, which also involve complex numbers, reinforces this supposition. The way in which the geometrical conception of imaginaries became familiar to him, and the results it could lead to in his hands, are clearly shown by the notes (published only recently) in which he applies complex numbers to the solution of problems of elementary geometry ({[}1{]}, vol.~4, p.~396 and vol.~8, p.~307). Even more explicit is the letter to Bessel in 1811 ({[}1{]}, vol.~8, p.~90-91) in which he sketches the essentials of the theory of integration of functions of a complex variable: ``Just as the whole domain of real quantities can be represented by means of an infinite line, so the complete domain of all quantities, both real and imaginary, can be realized by means of an infinite plane, in which each point, determined by its abscissa a and its ordinate b, represents at the same time the quantity \(a + ib\) . The continuous passage from one value of x to another consequently takes place along a curve, and can therefore be achieved in an infinity of ways\ldots''.

But it was not until 1831 that Gauss (à propos the introduction of ``Gaussian integers'' \(a + ib\) , where \(a\) and \(b\) are integers) publicly expounded his ideas on this point so precisely ({[}1{]}, vol.~2, Theoria Residuorum Biquadraticorum, Commentatis Secunda, art. 38, p.~109, and Anzeige, p.~174 et seq.). In the intervening years, the idea of representing complex numbers geometrically had been independently rediscovered by two modest amateurs, both of them more or less self-taught, who made no other contribution to mathematics, and who were both without much contact with the scientific circles of their time. For this reason their work was in danger of passing completely unnoticed; this is precisely what happened to the first in point of date, a pamphlet written by a Dane, C. Wessel. Published in 1798, it was clearly conceived and written, but was not rescued from oblivion until a century later. The same mishap might have overtaken the second, by the Swiss J. Argand, who owed it only to chance that he saw his work, published in 1806, exhumed seven years later (*). This work provoked an active discussion in the Annales de Gergonne, and was the subject, in France and England, of several publications by obscure authors between 1820 and 1830. But the authority of a great name was needed to put an end to these controversies and to rally mathematicians to the new point of view, and it was not until the middle of the century that the geometrical representation of complex numbers at last became universally adopted, following the publications of Gauss (mentioned above) in Germany, the work of Hamilton and Cayley on hypercomplex systems in England, and finally the support of Cauchy (** ) in France, only a few years before the genius of Riemann was to extend still further the role of geometry in the theory of analytic functions, and at the same stroke create the science of topology.

\par\noindent\textbf{* \(\ast\)}\par

The measurement of angles by means of the arcs they cut off on a circle is as old as the notion of angle itself, and was already known to the Babylonians: their unit of angular measure was the degree, which we still use. The Babylonians used only measures of angles lying between o and 360°; this was adequate for their purposes, which were primarily to fix the positions of celestial objects at determinate points of their apparent orbits and to construct tables of these orbits for scientific or astrological purposes.

Among the Greek geometers of the classical era, the notion of angle (Euclid's Elements, I, Definitions 8 and 9) was even more restricted, because it applied only to angles smaller than two right angles; and since on the other hand their theory of proportion and measurement was based on the comparison of arbitrarily large multiples of the magnitudes being measured, angles could not be measurable quantities for them, although of course they had the concepts of equal angles, of one angle being greater or less than another, and of the sum of two angles, provided that the sum did not exceed two right angles. Just as with the addition of fractions, the measurement of angles must have been in their eyes an empirical procedure, devoid of scientific value. This attitude is well illustrated by the admirable essay of Archimedes on spirals ({[}3{]}, vol.~2, pp.~1-121) in which, since he cannot define them by the proportionality of radius vector to angle, he gives a kinematic definition (Definition 1, p.~44; cf.~the statement of Proposition 12, p.~46) from which he succeeds in extracting, as the rest of the work shows, everything that the general notion of angular measure would have given him had he been in possession of it. As to the Greek astronomers, they seem to have been content to follow their Babylonian predecessors on this point as on many others.

Here too, as in the evolution of the concept of real number (cf.~the Historical Note on Chapter IV), the relaxation of the spirit of rigour during the decadence of Greek science brought a return to the ``naive'' point of view which, in some respects, comes nearer to our own than does the rigid Euclidean conception. Thus an ill-advised interpolator inserted the following famous proposition in Euclid (Euclid's Elements, VI, 33): ``Angles are proportional to the arcs which they cut out on a circle'' (*), and an anonymous scholar who comments on the ``proof'' of this proposition does not hesitate to introduce, of course without justification, arcs equal to arbitrarily large multiples of a circumference, and the angles corresponding to these arcs (**). But even Vietà, in the 16th century, although he appeared to come close to our modern conception of angle when he discovered that the equation \(\sin nx = \sin \alpha\) has several roots, obtained only those roots corresponding to angles smaller than two right angles ({[}4{]}, p.~305). Only in the 17th century was this point of view definitively superseded; and, after the discovery by Newton of the power-series expansions for \(\sin x\) and \(\cos x\) had furnished expressions of these functions which were valid for all values of the variable, Euler finally formulated the precise conception of the notion of angular measure, in connection with logarithms of ``imaginary'' numbers ({[}5{]}, (1), vol.~17, p.~220).

Of course, the classical definition of angular measure in terms of the length of a circular arc is not only intuitive but essentially correct; however it requires, to make it rigorous, the notion of the length of a curve, i.e., integral calculus. From the point of view of the structures which come into play, this is a very long-winded procedure, and it is possible, as we have seen in the text, to arrive at the same end by no other means than those of the theory of topological groups; in this manner the real exponential and the complex exponential appear as arising from the same source, the theorem which characterizes the ``one-parameter groups'' (Chapter V, § 3, Theorem 1).

